% !Mode:: "TeX:UTF-8" 

\BiChapter{实验设计与结果}{}

\BiSection{测试环境}{}

\BiSubsection{硬件环境}{}

\begin{itemize}
    \item CPU：AMD Ryzen 9 7950X 16-Core Processor @ 5.5GHz
    \item 内存：96GB DDR5 @ 4800MHz
    \item 硬盘：1TB NVMe SSD @ PCIe 3.0
\end{itemize}

\BiSubsection{软件环境}{}

\begin{itemize}
    \item 操作系统：Arch Linux
    \item 内核版本：Linux 6.14.6-arch1-1
    \item Java 版本: OpenJDK 24.0.1 2025-04-15
    \item Z3 版本: 4.13.4
\end{itemize}

\BiSection{测试指标}{}

\BiSubsection{难度}{}

为了更好地分析一个 IAM 策略的复杂程度，我们引入一个名为难度的度量。对一个 IAM 策略，我们约定：

\begin{enumerate}
    \item 对策略中的每一条语句，其难度值为1。
    \item 如果该语句包含一个通配符 (wildcard)，则其难度值为2。
    \item 如果该语句包含 $n$ 个通配符 (wildcard)，则其难度值为 $2^n$。
    \item 这个策略的难度值为所有语句的难度值的加和。
\end{enumerate}

例如对于策略\ref{policy_file_example}，其难度值为 10。

\BiSubsection{正确性}{}

以 AWS IAM 云服务给出的分析结果为参照，本地形式化验证系统的结果应该与云服务的结果相同。

\BiSubsection{性能}{}

分别记录 AWS IAM 云服务和本地系统在同一组难度递进的策略集上的求解用时，二者应该展现出相似的性能变化趋势。

\BiSection{实验设计}{}

对本地的性能测试，只需直接在代码中加入计时器即可，而对 AWS IAM 服务等计时则要棘手一些。AWS IAM 服务提供了命令行工具 AWS CLI，可以让用户在命令行中上传策略文件，执行分析并得到结果。具体涉及到的命令有：

\begin{lstlisting}[language=Bash, caption={AWS CLI 命令示例}, label={aws_cli}]
aws s3api put-bucket-policy --bucket <bucket> --policy file://<file_name>
aws accessanalyzer start-resource-scan \
    --analyzer-arn <arn> --resource-arn <arn>
aws accessanalyzer get-analyzed-resource \
    --analyzer-arn <arn> --resource-arn <arn>
aws accessanalyzer list-findings --analyzer-arn <arn>
\end{lstlisting}

其中，第一条命令用于将本地的策略文件上传到 AWS IAM 云服务器，第二条命令用于启动分析流程，第三条命令用于获取分析情况，包括最终状态和分析用时等，最后一条命令用于获取分析结果，即所有的发现项。然而，根据 AWS 文档，启动分析流程等命令是异步的，即本地命令行执行该命令后立即结束，而非等待云端等分析结果完成。所以为了获取分析用时，需要解决这个异步的情况。

采用轮询的机制来处理异步分析的情况。具体而言，当执行完启动分析命令后，开始周期性地执行获取分析情况的命令（如一秒一次），并对获取到的分析情况进行解析，当其中的状态量为 ACTIVE 时，代表分析已经完成，此时读取其分析用时，即可正确地对 AWS IAM 云服务进行计时。

在本地系统上运行测试时，将运行两次。第一次为预热，保证缓存已经被预填充，以避免可能的性能损耗，第二次为正式运行，将记录系统的运行时间。

\begin{algorithm}[H]
    \caption{GetFinding \label{get_finding}}
    \IncMargin{2em}
    \DontPrintSemicolon

    isActive = $False$ \tcp*{初始化分析状态}
    \While{$\lnot \text{isActive}$}{
        current $\gets$ \text{get-analyzed-resource} \tcp*{获取当前分析情况}
        isActive $\gets$ current.status \tcp*{更新分析状态} 
    }
\end{algorithm}

\BiSection{实验结果}{}

\BiSubsection{正确性}{}

经验证，本地形式化验证系统得出的分析结果和 AWS IAM 云服务得出的完全相同，证明本系统的实现是正确的。

\BiSubsection{性能}{}

随机构建的 15 个策略的运行情况如下：
\begin{table}[H]
	\renewcommand{\arraystretch}{1.2}
	\centering\wuhao
	\caption{求解性能量表} \label{tab_ch2} \vspace{2mm}
	\begin{tabularx}{\textwidth} { 
   >{\centering\arraybackslash}X 
   >{\centering\arraybackslash}X 
   >{\centering\arraybackslash}X 
   >{\centering\arraybackslash}X }
	\toprule[1.5pt]
		序号 & 难度 & AWS IAM 云用时(s) & 本地系统用时(ms) \\
	\midrule[1pt]
		1 & 8 & 11 & 31 \\
		2 & 16 & 17 & 425 \\
		3 & 24 & 19 & 1908 \\
        4 & 32 & 29 & 2906 \\
        5 & 40 & 63 & 6562 \\
        6 & 48 & 101 & 11122 \\
        7 & 56 & 159 & 18364 \\
        8 & 64 & 245 & 27498 \\
        9 & 72 & 393 & 46155 \\
        10 & 80 & 567 & 69846 \\
        11 & 88 & 769 & 内存不足 \\
        12-15 & 96-120 & 求解失败 & 内存不足 \\
	\bottomrule[1.5pt]
	\end{tabularx}
\end{table}

绘制性能变化趋势图如下：

\begin{figure}[H]
\centering
\includegraphics[width=\textwidth]{scalability05}
\caption{求解性能变化趋势图} \label{performance}
\end{figure}

观察发现，本地系统求解性能变化趋势与 AWS IAM 云服务求解性能变化趋势基本一致，但是本地系统变化趋势的常系数项更大，这是因为本地系统的硬件性能不如 AWS IAM 云计算中心导致的。